\section{Passage to the radial series}\label{c4:radial}

The angular identities are now established with their coefficientwise meaning and convergence domain.  We next transform the oscillator and momentum factors.  The resulting radial series must retain its phases, its invariant momentum line, and the order of the period and compensator limits.


We continue the oscillator notation, including cases with only screenings. Set \(R=A,\ \Pi=p(p,\cdot)/p^2,\ K_R=(I-\Pi)K(I-\Pi)\), which acts as \((R-I)^{-1}\) on \(p^\perp\) by the solution of \((I-R)m=d\) above. On the invariant line use positive frequencies \(l\in\mathbb Z_{>0}\), operators \(\Lambda_l=l I\) there with projection \(P_l=\Pi\). On a noninvariant generalized eigenblock, of eigenvalue \(r=\exp(2\pi i\gamma)\), \(0<\gamma<1\) (\(\gamma=1/6,1/3,2/3,5/6\)), put \(\Gamma=(2\pi i)^{-1}\log R\) on that branch. Positive frequencies here are \(l=n'+\gamma,\ n'\ge0\) integral, with operators \(\Lambda_l=n'+\Gamma\) on this block and \(P_l\) its spectral projection. Negative frequencies, when used, take the same definitions with \(n'<0\) (or negative integers on the invariant line), so \(P_l^*=P_{-l},\ (\Lambda_l P_l)^*=-\Lambda_{-l}P_{-l}\). All matrix functions are ordinary finite-dimensional holomorphic calculus on these blocks, using at most derivatives for the nilpotent parts. Write
\[
 H_0(z)=\sum_{n\ge0}\left(\frac1{n+1}-\frac1{n+z}\right),
\]
and denote by \(H_0[R]\) the operator using \(H_0(\Gamma)\) on noninvariant blocks and zero on the invariant line.

Here is the radial version of \eqref{c4:Ang-2}. Set
\[
 Z_R^r(L)=\prod_{l>0}\det(1-e^{-L\Lambda_l})^{-1},\quad
 Z_C^r(L)=\prod_{\substack{n\ge1\\8\nmid n}}(1-e^{-nL/8})^{-1},\quad
 \delta=(\mu,p),\quad a=\delta/p^2
\]
(\(a=7/18\), or \(7/18+1\) for the shifted baseline).
There is a prefactor \(\exp(7L/192) Z_R^r(L)/Z_C^r(L)\). Use the same screening expansion and coefficients/measures as there, now sorting charges by increasing \(u=\Re s\), periodically repeating them with period \(L\) on the real axis (same \(\phi\)). We can use any base interval of length \(L\) containing the deterministic vertices at 0 and \(s_*\) when present (keeping them unshifted; screening powers in the angular argument are integral). Ignore ties, a null set. Write
\[
 n_b=(p,b),\qquad \zeta_b=(p,K b).
\]
Keep \(\sum_b n_b=0\) per period. Sum over a slope initially \(\ell_0\in a+\mathbb Z/p^2\) at the left end, stepping by \(n_b/p^2\) at each vertex of charge \(b\) (traversing rightward); \(\ell_b\) denotes its value midway through the step. Use weight \(\exp(-\int p^2\ell(u)^2du/2)\) over the period and, per vertex, the factor
\begin{equation}
 \exp\!\big((H_0[R]b,b)/2 + 2\pi i(\mu,K_R^*b)+2\pi i\zeta_b(\ell_b-a)+i\phi_b n_b\ell_b\big).             \label{c4:Rad-1}
\end{equation}
In the marked case multiply further by \(\exp(2\pi i(b_*,Kp))\) for compensator charge \(b_*\), and by \(\exp(-2\pi i\zeta_b)\) at each screening strictly between 0 and \(u_*\). Finally each vertex \(x\) in the base period and \(y\) anywhere strictly later on the periodic extension give a factor \(\exp(-\mathcal Q(x,y))\),
\begin{equation}
 \mathcal Q(x,y)=\sum_{l>0}
     \big(\Lambda_l^{-1}\exp(-\Lambda_l(s_y-s_x))P_l b_x,b_y\big) .                       \label{c4:Rad-2}
\end{equation}
Here \(s\) includes the real period shift. This counts also later self images.

\subsection{Fourier transformation and the momentum sum}

We derive the formula. Put \(\kappa=4\pi^2/L\). For \(j<k\) in angular order let \(d'=s_j-s_k=u+i\phi,\ -2\pi<\phi<0\) initially. The oscillator pair logarithm uses the matrix function \(G_R(d')\), paired as \((G_R(d')b_j,b_k)\), where
\[
 G_r(d')=\sum_{m\ge0}r^m\log(1-e^{-\kappa(m+i d'/(2\pi))})
          +\sum_{m<0}r^m\log(1-e^{\kappa(m+i d'/(2\pi))}).
\]
For scalar \(r=e^{2\pi i\gamma}\) (also \(\gamma=0\) for the invariant) let \(l\) range over \(\gamma+\mathbb Z\), excluding 0. Define for \(-L<u<0\), extended periodically in \(u\),
\[
 \mathcal H_r(d')=-\sum_{l>0}\frac{e^{l d'}}{l(1-e^{-lL})}
                -\sum_{l<0}\frac{e^{l(d'+L)}}{|l|(1-e^{-|l|L})}.
\]
Near \(0<\gamma<1\) use analytic branches (thus \(|l|=-l\) for the negative frequencies). For \(-L<u<L,\ u\ne0\), one has
\[
 \begin{aligned}
 G_r&=\mathcal H_r + (2\pi i/L)(r-1)^{-1}d'
           -\pi i\,\operatorname{sgn}(u)/(r-1)-\kappa r/(r-1)^2,\quad r\ne1,\\
 G_1&=\mathcal H_1-\big((d')^2/2+i\pi d'-\pi^2/3\big)/L
                     -L/12+(d'+i\pi)\operatorname{sgn}(u)/2.
 \end{aligned}
\]
Indeed the Fourier coefficients in \(u\), mode \(\nu=2\pi k/L\), of \(G_r\) are \(-e^{-\nu\phi}/[k(r e^{2\pi\nu}-1)]\) for \(k\ne0\), 0 at \(k=0\). For \(\mathcal H_r\) they are \(-\sum_l e^{il\phi}/[L l(l-i\nu)]\) by integrating each half separately. Including zero where applicable,
\[
 \sum_{l\in\gamma+\mathbb Z} \frac{e^{il\phi}}{l-i\nu}
 =\mathcal J(\nu)=\frac{2\pi i\,e^{-\nu\phi}}{r e^{2\pi\nu}-1}
\]
(the Fourier series of this twisted exponential on \(-2\pi<\phi<0\)). Thus for \(r\ne1,\ k\ne0\), the coefficient of \(G-\mathcal H\) is \(-\mathcal J(0)/(L i\nu)\), and its constant coefficient is \(\mathcal J'(0)/(iL)\), matching the stated corrections. At \(r=1\) exclude the zero term \(i/\nu\) in the \(\mathcal H\) calculation, using
\(\mathcal J(\nu)=i/\nu-i(\phi+\pi)+i\nu(\phi^2/2+\pi\phi+\pi^2/3)+O(\nu^2)\); this gives coefficients \(((\phi+\pi)/\nu-1/\nu^2)/L\) and \((\phi^2/2+\pi\phi+\pi^2/3)/L\) respectively, just those of the second correction. Take limits also to \(\phi=0^-\) for same-contour interactions. Substitute the blocks (where \(r/(r-1)^2\) gives \(-K_R^*K_R\)). The \(\mathcal H\) part sums precisely minus \eqref{c4:Rad-2} over the images for these two vertices, using bilinear adjunction to express both radial orders.

For the self product including scaling take half the log pair with itself, \(d'\to0\) in \(u<0,\phi<0\), subtracting \(b^2\log(1-e^{-2\pi i d'/L})\) but adding \(b^2\log(2\pi/L)\) before halving. Here \(\mathcal H_r-\log(-d')\) tends to
\[
 H_0(\gamma)-\sum_l\frac{e^{-|l|L}}{|l|(1-e^{-|l|L})}
\]
(replace \(H_0(\gamma)\) by 0 at \(r=1\)); the unwrapped positive-frequency sum differs in the limit from \(\log(1-e^{d'})\) by the defining series of \(H_0\). This is also analytic near noninteger \(\gamma\), e.g. by differentiating in \(d'\) using the geometric sum. Difference of the singular log and the scaled subtraction contributes \(i\pi b^2/2\) before halving, canceling the imaginary sign constants since \(K_R+K_R^*=-(I-\Pi)\). Thus the self log gives the \(H_0[R]\) term of \eqref{c4:Rad-1}, minus \eqref{c4:Rad-2} for self images, plus
\[
 [\kappa(K_R b)^2+(\pi^2/(3L)-L/12)(\Pi b)^2]/2 .
\]

The state momentum \(n\) of \eqref{c4:Ang-2} is \(n=K_R^*d+M p\), \(M=t-Z/p^2,\ Z=\sum\zeta_b,\ t\in\mathbb Z,\ d=\sum b\), \((p,d)=0\).
The order-independent linear-in-\(d'\) terms just computed, operator \((2\pi i/L)(K_R-\Pi/2)\), cancel the angular \(z\)-argument powers except for \(-2\pi i M S'/L,\ S'=\sum_b n_b s_b\). Indeed the coefficient at \(s_j\) groups using \((K_R b_j,d)+(b_j,\sum_{k<j}b_k+b_j/2)\). The noninvariant \(\kappa\) corrections cancel the corresponding angular heat; invariant constants cancel by \(\Pi d=0\). The \((d')^2\) terms give \((S')^2/(2Lp^2)\). Together the \(M\) sum (including Cartan phase) now uses
\[
 \exp\!\big(-\kappa p^2(M+i S'/(2\pi p^2))^2/2+2\pi i\delta M\big).
\]
Poisson summation (Fourier expansion of this periodized Gaussian) gives the scalar \(\sqrt{L/(2\pi p^2)}\) and sum over the stated \(\ell_0\) weighted by
\[
 \exp(-p^2 L\ell_0^2/2+\ell_0 S'+2\pi i(\ell_0-a)Z).
\]
Adding invariant \(d'\operatorname{sgn}(u)\) pair terms gives the heat integral and \(i\phi_b n_b\ell_b\) phases: the linear-in-\(s\) contribution is \(\sum_b s_b n_b[\ell_0+(N_{<b}+n_b/2)/p^2]\), where \(N_{<b}\) is total \(n_c\) earlier radially, so real-position contributions telescope with the slope squares.

Remaining pair phases from cocycles and \(G-\mathcal H\) use, dividing the log by \(\pi i\), the angular \(j<k\) sum of
\[
 (1-\operatorname{sgn}u)(K b_j,b_k)+\operatorname{sgn}(u)(U b_j,b_k)+2(b_j,K\sigma_k),\qquad U=K-K_R+\Pi/2.
\]
Here \((U b_j,b_k)=(\zeta_{b_j} n_{b_k}-n_{b_j}\zeta_{b_k})/p^2\), so the signed \(U\) sum is \(\sum_b(2N_{<b}+n_b)\zeta_b/p^2\), giving the \(\zeta_b\) term in \eqref{c4:Rad-1} together with the Poisson phase. In the rest, all roots precede external charges angularly and \((K b_j,b_k-\sigma_k)\) is integral for root \(b_j\). Thus modulo 2 only the marked case matters: for such \(j\) only external \(k\) to its radial left leave a \(2(b_j,K\sigma_k)\) phase term, congruent to \(-2\zeta_{b_j}\) from \(B\) and \(+2\zeta_{b_j}\) from \(b_*\). The external pair itself contributes \(2(b_*,Kp)\) since \(u_*>0\). This is the claimed rule.

For the prefactor put \(\mathfrak p(y)=\prod_{m\ge1}(1-y^m)\). The scalar invariant self calculation above (both charges 1 on a Euclidean unit line) gives
\[
 \mathfrak p(e^{-\kappa})=\sqrt{L/(2\pi)}\,\exp((\kappa-L)/24)\mathfrak p(e^{-L})
\]
using \(-\sum_{l\ge1} e^{-lL}/[l(1-e^{-lL})]=\log\mathfrak p(e^{-L})\). Characteristic polynomials give \(Z_C=\mathfrak p(Q_L)/\mathfrak p(Q_L^8)\), \(Z_R=\mathfrak p(Q_L^2)/[\mathfrak p(Q_L^3)\mathfrak p(Q_L^6)]\). Apply the identity there. \(Z_C\) becomes \(\sqrt{8}\exp(-7\kappa/24-7L/192)Z_C^r\), and \(Z_R\) together with the Poisson scalar becomes \(\sqrt{8}\exp(-7\kappa/24)Z_R^r\) where 
\[Z_R^r=\mathfrak p(e^{-L/2})/[\mathfrak p(e^{-L/3})\mathfrak p(e^{-L/6})].\]
This proves the prefactor.

\subsection{The marked interval and its root activities}

The transformed phases now determine which momentum lattice applies between the marks. The large central charge will suppress all but one type of root in that interval.

Some useful consequences:
\par\noindent\textbullet\quad In the marked interval \((0,u_*)\) the slope lattice is \(a+1+\mathbb Z/p^2\) since \(n_B=p^2\); the extra interval character there makes the single-root factors \eqref{c4:Rad-1} agree with the shifted baseline (\(a\) replaced by \(a+1\)). Elsewhere they use the ordinary baseline.
\par\noindent\textbullet\quad Since \(h\) is on the frequency \(\pm1/6\bmod1\) plane of \(R\), the \(Nh\) part of \(B\) gives the self factor \(d_R(h)^{N^2}\) with
\[
 d_R(h)=6^{-h^2/2}\prod_{m=1}^5|1-e^{2\pi i m/6}|^{(R^m h,h)/2}.
\]
Indeed \(d_R(h)=\exp((H_0[R]h,h)/2)\), by writing the unwrapped sum from the \(\mathcal H\) self calculation as \(\sum_{m=0}^5(R^m h,h)\log(1-\exp((d'-2\pi i m)/6))\) and taking \(d'\to0^-\). Also \(\zeta_h=-9/8=-p^2\) by the above \(Kh\) computation; at \(B\) the pre-step slope minus \(a\) lies in \(\mathbb Z/p^2\). Thus all local phases from \(Nh\) are global unit scalars independent of the configuration and compensator summand.
\par\noindent\textbullet\quad Removing \(Nh\) from \eqref{c4:Rad-2}, besides its interactions with images and compensator, multiplies the activities at root vertices by \(\exp(N F_b(s))\), periodically imaged. On the infinite axis with just the \(B\) at 0,
\[
 \Re F_b(s)=\sum_{m=0}^5(h,R^m b)\log|1-\exp((-|u|+i\phi-2\pi i m)/6)|.
\]
This follows by the positive-frequency sum as just used (either radial order, using adjunction). For columns \eqref{c4:Ang-1} these rows of exponents as \(m\) varies are respectively \(x',x'+y',y',x',x'+y',0,y'\), with \(x'=(0,-1,-1,0,1,1),\ y'=(1,1,0,-1,-1,0)\). For each nontrivial pattern the sum is \(\le-c e^{-|u|/6}\) uniformly on the chosen root rays close enough to \(-\pi\). Indeed at \(-\pi\) each of the \(x',y'\) sums is the log ratio of distances at cosines \(\sqrt{3}/2\) and \(-\sqrt{3}/2\), and continuity preserves the estimate (also of first derivative at zero radius). The neutral type \(\beta\) gives \(F\equiv0\), and \(n_\beta=1\). Each other
\[
 F_b(u+i\phi)=c_b^\pm e^{-|u|/6}+O(e^{-5|u|/6}),\qquad \Re c_b^\pm<0,\qquad u\to\pm\infty ,
\]
by the frequencies coupling to \(h\).

\subsection{A coercive bound for the radial interaction}
\label{c4:radial-coercivity}

We give the positivity estimate that justifies convergence despite the indefinite form. Fix the finite root rays for the factors of \eqref{c4:Ang-1}, always keeping distinct angles for distinct integration factors. Their spread around \(-\pi\) can be chosen small enough so that for any finite root configuration on the real line (distinct real positions) and weights \(t_x\ge0\),
\begin{equation}
 \Re\sum_{x<y}t_x t_y\mathcal Q(x,y)\ \ge\
       c\sum_I\left(\sum_{x\in I}t_x\right)^2-C\sum_x t_x^2,                     \label{c4:Rad-3}
\end{equation}
\(I\) an ordinary translated unit cell partition, with constants independent of positions and partition.

Here are the form matrices. Put \(T=(R+R^{-1})/2\), \(B_0\) the seven columns of \eqref{c4:Ang-1}. Write
\[
 B_0^*(X-T)^{-1}B_0
 = M_0/(X-1)+M_1/(X-1/2)+M_2/(X+1/2)+M_3/(X+1/2)^2
\]
(\(B_0^*\) pairs against the form). \(M_0,M_1,M_3\) are positive semidefinite, \((M_1)_{jk}=1/2\) on \(j,k\in\{1,4\}\). Also \(M_2=Y Y^t/Y_1+E\) with \(Y=(22,9,27,-22,-18,9,0)^t/36\) and \(E\) entrywise nonnegative with strictly positive diagonal outside \(\{1,4\}\). For verification \(M_0\) is the invariant \(p\)-Gram. \(M_i=C_i^t D_i C_i\) for \(i=1,3\), with
\[
 \begin{split}
 D_1=\begin{pmatrix}2&1\\1&2\end{pmatrix}/4,&\quad
 C_1=\begin{pmatrix}1&0&-1&1&0&0&-1\\0&1&1&0&1&0&1\end{pmatrix},\\
 D_3=\begin{pmatrix}2&1\\1&2\end{pmatrix}/12,&\quad
 C_3=\begin{pmatrix}1&0&1&-1&-2&0&-1\\0&1&1&0&1&2&1\end{pmatrix}.
 \end{split}
\]
And
\[
 36M_2=\begin{pmatrix}
22&23&45&-22&-9&40&13\\23&22&27&-5&18&32&23\\45&27&54&-27&-9&36&18\\
-22&-5&-27&22&27&-4&5\\-9&18&-9&27&54&36&27\\40&32&36&-4&36&40&40\\13&23&18&5&27&40&22
 \end{pmatrix}.
\]
These result by multiplying \(B_0^*,B_0\) against \(P_+=-2(T-I)(T+I/2)^2,\ P_-=I-\Pi-P_+,\ (T+I/2)P_-\) for \(M_1,M_2,M_3\) respectively.

Smooth each oriented separation in \eqref{c4:Rad-2} by adding a small \(\eta>0\), and temporarily include diagonals with half weight. The resulting real kernel in the form is, up to entrywise error \(O((\eta+\eta_1)(1+u^3)e^{-u/6})\), \(\eta_1=\max_b|\phi_b+\pi|\), \(u\ge0\) the absolute real separation and \(\phi\) angle difference,
\[
 M_0 K_{\rm inv}^\eta(u,\phi)+M_1 K_{1/2}^\eta(u,\phi)+M_2 K_{-1/2}^\eta(u,\phi)
       +M_3\,\partial_X K_X^0(u,0)|_{X=-1/2},
\]
where \(K_{\rm inv}^\eta=\sum_{l\ge1\ {\rm integer}} e^{-l(u+\eta)}\cos(l\phi)/l\) and
\[
 K_X^\eta(u,\phi)=\frac12\sum_{\substack{l>0\\\cos(2\pi l)=X}} e^{-l(u+\eta)}\cos(l\phi)/l .
\]
Indeed evaluate \(f_\gamma(D')=\sum_{n\ge0}e^{-(n+\gamma)D'}/(n+\gamma)\) as analytic function on the monodromy branches, \(D'=u+\eta+i\phi\), with \(\gamma=1\) on the invariant line. \(f_\gamma(D')+\log(1-e^{-D'})\) is regular at \(D'=0\) also under \(\gamma\)-differentiation by the geometric-sum derivative. Since \(R\) is real and charges real, taking the real part replaces \(f_\gamma(D')\) on noninvariant blocks by the branch function \((f_\gamma(D')+f_{1-\gamma}(\overline{D'}))/2\). Further symmetrizing under \(\gamma\leftrightarrow 1-\gamma\) costs only the stated smooth angle error by the common singular term (at large \(u\) use the exponential sums directly including nilpotent derivatives). This gives \(K_X^\eta\) at \(X=\cos(2\pi\gamma)\), locally analytic in \(X\), hence the form matrices of \(T\); the derivative term may likewise be taken at \(\eta=\phi=0\) by the same estimates.

\(K_{\rm inv}^\eta,K_{\pm1/2}^\eta\) and the indicated derivative (on \(u\) differences with angle ignored) are all positive semidefinite kernels. For the former this is termwise from \(e^{-l|u|}\) and \(\cos(l\phi)\) on differences. For the derivative use the real-line Fourier transform of \(K_X^0\) at \(\phi=0\):
\[
 \sum_{l\in\gamma+\mathbb Z}\frac1{l^2+\xi^2}
     =\frac{\pi}{\xi}\frac{\sinh(2\pi\xi)}{\cosh(2\pi\xi)-X}.
\]
This follows, for instance, from the same \(\mathcal J\) formula for both signs at \(\phi\to0^-\); its \(X\)-derivative is positive with the continuous limit at \(\xi=0\). Kernel differentiation here has no log singularity as checked above, and inversion gives the assertion. Also \(K_{-1/2}^\eta\ge c e^{-u/3}\) pointwise with small angles and smoothing (first term for large \(u\), positive logarithmic singularity near zero, continuity and strict positivity at zero angle in between).

Thus after discarding positive semidefinite pieces (matrix Gram times scalar kernel), \(E\) controls with a positive constant the squared unit-cell weights for types outside \(\{1,4\}\). For \(\{1,4\}\) retain also a small fraction of the \(l=1/6\) term with \(M_1\); restricted to these indices it controls squared unit-cell weights there by the pointwise bound, and on a general vector still gives half that bound minus a constant times the squared cell weights for remaining types, by the squared norm inequality and kernel decay. Choose the retained fraction small enough to absorb this loss into the \(E\) bound. This leaves \(c>0\) for total squared cell weights, uniform for small spreads and \(\eta\). The errors estimated above are absorbed by taking both small (a kernel bounded by a decreasing exponential at infinity and bounded at zero costs at most a constant times total squared cell weights by cell convolution). Now remove the smoothed diagonals at cost \(C(\eta)\sum t_x^2\). Original distinct-ray kernels differ by \(O(\eta)\) with similar decay for fixed positive angle gaps. On like-ray pairs the leading \(-2\log|1-e^{-u}|\) singularity for the identical root gives only favorable change from smoothing, and smooth remainders again differ by \(O(\eta)\) with decay. Taking \(\eta\) sufficiently small also relative to angle gaps now proves \eqref{c4:Rad-3}.

For periodic repetition apply \eqref{c4:Rad-3} to many periods on the line, divide by period count and pass to its limit (interactions decay over large separations). This gives in particular \(c k_{\rm tot}^2/L-C k_{\rm tot}\) lower bound on real root interaction per period by cell Cauchy--Schwarz, \(k_{\rm tot}\) the number of roots. Under a small complex change of \(L\) scaling all \(u\) proportionally, the kernel perturbations are uniformly small with exponential decay in original real separation, including for identical roots arbitrarily close by the common log singularity and bounded log-ratio estimate. For each fixed reference \(L\) their image sums are absorbed into the quadratic gap. External vertices on separated rays cost only linearly in \(k_{\rm tot}\); phases remain unitary times constant activities, and the absolute sum over \(\ell_0\) of the heat cost is \(O_L(k_{\rm tot}+1)\) since slope differs from entrance by \(O(k_{\rm tot}+1)\). Hence the resulting full series converges absolutely locally uniformly and is holomorphic near each real \(L>0\) by integration on scaled coordinates with fixed real ordering. For fixed numbers of vertices all transformations above are exact after the heat sum, so by the absolute estimates on both sides near small \(Q_L\) this indeed gives the continued value of \eqref{c4:Ang-2} and both baseline identities.

